% GENERATED FILE. Edit canonical lessons, then run npm run generate:books.
% canonical-source-hash: fnv1a64:3d105b3b601989dd
% canonical-lessons: IT-C377-disponibile, IT-C377-impegnato, IT-C377-impegno, IT-C377-impegnarsi, IT-C377-flessibile

\chapter{Se ho tempo}
\label{ch:it-se-ho-tempo}
\hlchaptermodality{377}

\begin{chapteropening}
\textbf{By the end of this chapter:} \emph{I can say whether I am free or busy, talk about a commitment and say I am flexible.}

Complete the last lesson of chapter 377: I can say whether I am free or busy, talk about a commitment and say I am flexible.
\end{chapteropening}

\section[disponibile]{disponibile — available; willing, helpful}

\label{lesson:IT-C377-disponibile}



\begin{quote}
Before the new one: say the Italian for unlikely, then the Italian for the probability.
\end{quote}

\subsection*{You'll want to know: disponibile}
\textbf{disponibile} — \textquotedblleft{}available; willing, helpful\textquotedblright{}.

\begin{quote}
Se sei disponibile sabato, ci vediamo.
\end{quote}

\begin{quote}
\emph{If you are free on Saturday, let us meet.}
\end{quote}

Of a person, also willing to help: \emph{è sempre molto disponibile}.

\begin{cousinweb}
\emph{disponibile} is built on \emph{disporre} \textquotedblleft{}to arrange, to have at one's disposal\textquotedblright{}, from Latin \emph{dispōnere} \textquotedblleft{}to set in order\textquotedblright{} (Lewis \& Short), built on \emph{pōnere} \textquotedblleft{}to put\textquotedblright{}.
\end{cousinweb}

\begin{grammarlens}[title={What you've built}]
One more word for telling, explaining and describing at B1.
\end{grammarlens}

\subsection*{Your turn}
\begin{itemize}
\raggedright
  \item \emph{Say it:} \emph{disponibile}
  \item \emph{Say it:} \emph{disponibile}, once more
  \item \emph{Say it:} a sentence of your own with \emph{disponibile}
  \item \emph{Recall:} say the Italian for unlikely, then the Italian for the probability, then say \emph{disponibile} again
\end{itemize}

\begin{tcolorbox}[breakable,colback=teal!4,colframe=teal!35!black,title={Before you move on}]
What does disponibile mean? (\textquotedblleft{}Available; willing, helpful\textquotedblright{}.) Say it once more.
\end{tcolorbox}
\section[impegnato]{impegnato — busy, tied up; committed}

\label{lesson:IT-C377-impegnato}



\begin{quote}
Before the new one: say the Italian for the probability, then the Italian for available.
\end{quote}

\subsection*{You'll want to know: impegnato}
\textbf{impegnato} — \textquotedblleft{}busy, tied up; committed\textquotedblright{}.

\begin{quote}
Domani sono impegnato tutto il giorno.
\end{quote}

\begin{quote}
\emph{Tomorrow I am busy all day.}
\end{quote}

Also socially engaged: \emph{un cantante impegnato}.

\begin{cousinweb}
\emph{impegnato} is the past participle of \emph{impegnare} \textquotedblleft{}to pledge, to engage\textquotedblright{}, built from \emph{in-} and \emph{pegno} \textquotedblleft{}a pledge\textquotedblright{}, inherited from Latin \emph{pignus} \textquotedblleft{}a pledge\textquotedblright{} (Lewis \& Short).
\end{cousinweb}

\begin{grammarlens}[title={What you've built}]
One more word for telling, explaining and describing at B1.
\end{grammarlens}

\subsection*{Your turn}
\begin{itemize}
\raggedright
  \item \emph{Say it:} \emph{impegnato}
  \item \emph{Say it:} \emph{impegnato}, once more
  \item \emph{Say it:} a sentence of your own with \emph{impegnato}
  \item \emph{Recall:} say the Italian for the probability, then the Italian for available, then say \emph{impegnato} again
\end{itemize}

\begin{tcolorbox}[breakable,colback=teal!4,colframe=teal!35!black,title={Before you move on}]
What does impegnato mean? (\textquotedblleft{}Busy, tied up; committed\textquotedblright{}.) Say it once more.
\end{tcolorbox}
\section[l'impegno]{l'impegno — the commitment, the engagement; the effort}

\label{lesson:IT-C377-impegno}



\begin{quote}
Before the new one: say the Italian for available, then the Italian for busy.
\end{quote}

\subsection*{You'll want to know: l'impegno}
\textbf{l'impegno} — \textquotedblleft{}the commitment, the engagement; the effort\textquotedblright{}.

\begin{quote}
Ho un impegno, ma se finisco presto ti chiamo.
\end{quote}

\begin{quote}
\emph{I have something on, but if I finish early I will call you.}
\end{quote}

\emph{Con impegno} — with effort, seriously.

\begin{cousinweb}
\emph{impegno} is built on \emph{impegnare} \textquotedblleft{}to pledge\textquotedblright{}, from \emph{pegno} \textquotedblleft{}a pledge\textquotedblright{}, inherited from Latin \emph{pignus} \textquotedblleft{}a pledge, a security\textquotedblright{} (Lewis \& Short).
\end{cousinweb}

\begin{grammarlens}[title={What you've built}]
One more word for telling, explaining and describing at B1.
\end{grammarlens}

\subsection*{Your turn}
\begin{itemize}
\raggedright
  \item \emph{Say it:} \emph{l'impegno}
  \item \emph{Say it:} \emph{l'impegno}, once more
  \item \emph{Say it:} a sentence of your own with \emph{l'impegno}
  \item \emph{Recall:} say the Italian for available, then the Italian for busy, then say \emph{l'impegno} again
\end{itemize}

\begin{tcolorbox}[breakable,colback=teal!4,colframe=teal!35!black,title={Before you move on}]
What does l'impegno mean? (\textquotedblleft{}The commitment, the engagement; the effort\textquotedblright{}.) Say it once more.
\end{tcolorbox}
\section[impegnarsi]{impegnarsi — to commit yourself, to make an effort}

\label{lesson:IT-C377-impegnarsi}



\begin{quote}
Before the new one: say the Italian for busy, then the Italian for the commitment.
\end{quote}

\subsection*{You'll want to know: impegnarsi}
\textbf{impegnarsi} — \textquotedblleft{}to commit yourself, to make an effort\textquotedblright{}.

\begin{quote}
Se ti impegni, passi l'esame.
\end{quote}

\begin{quote}
\emph{If you make an effort, you will pass the exam.}
\end{quote}

\emph{Impegnarsi a fare qualcosa} — to undertake to do something.

\begin{cousinweb}
\emph{impegnarsi} is built on \emph{impegnare} \textquotedblleft{}to pledge\textquotedblright{}, from \emph{pegno} \textquotedblleft{}a pledge\textquotedblright{}, inherited from Latin \emph{pignus} (Lewis \& Short): to give yourself as a pledge.
\end{cousinweb}

\begin{grammarlens}[title={What you've built}]
One more word for telling, explaining and describing at B1.
\end{grammarlens}

\subsection*{Your turn}
\begin{itemize}
\raggedright
  \item \emph{Say it:} \emph{impegnarsi}
  \item \emph{Say it:} \emph{impegnarsi}, once more
  \item \emph{Say it:} a sentence of your own with \emph{impegnarsi}
  \item \emph{Recall:} say the Italian for busy, then the Italian for the commitment, then say \emph{impegnarsi} again
\end{itemize}

\begin{tcolorbox}[breakable,colback=teal!4,colframe=teal!35!black,title={Before you move on}]
What does impegnarsi mean? (\textquotedblleft{}To commit yourself, to make an effort\textquotedblright{}.) Say it once more.
\end{tcolorbox}
\section[flessibile]{flessibile — flexible}

\label{lesson:IT-C377-flessibile}



\begin{quote}
Before the new one: say the Italian for the commitment, then the Italian for to commit yourself.
\end{quote}

\subsection*{You'll want to know: flessibile}
\textbf{flessibile} — \textquotedblleft{}flexible\textquotedblright{}.

\begin{quote}
Ho un orario flessibile: posso cominciare alle dieci.
\end{quote}

\begin{quote}
\emph{I have flexible hours: I can start at ten.}
\end{quote}

Of things too: \emph{un tubo flessibile} — a hose.

\begin{cousinweb}
\emph{flessibile} is a learned borrowing from Latin \emph{flexibilis} \textquotedblleft{}pliant, easily bent\textquotedblright{} (Lewis \& Short), from \emph{flectere} \textquotedblleft{}to bend\textquotedblright{}.
\end{cousinweb}

\begin{grammarlens}[title={What you've built}]
One more word for telling, explaining and describing at B1.
\end{grammarlens}

\subsection*{Your turn}
\begin{itemize}
\raggedright
  \item \emph{Say it:} \emph{flessibile}
  \item \emph{Say it:} \emph{flessibile}, once more
  \item \emph{Say it:} a sentence of your own with \emph{flessibile}
  \item \emph{Recall:} say the Italian for the commitment, then the Italian for to commit yourself, then say \emph{flessibile} again
\end{itemize}

\begin{tcolorbox}[breakable,colback=teal!4,colframe=teal!35!black,title={Before you move on}]
What does flessibile mean? (\textquotedblleft{}Flexible\textquotedblright{}.) Say it once more.
\end{tcolorbox}
